\documentclass[11pt,a4paper]{article}

% ====================== PAKIETY ======================
\usepackage[utf8]{inputenc}
\usepackage[T1]{fontenc}
\usepackage[polish]{babel}
\usepackage[a4paper,margin=2.3cm]{geometry}
\usepackage{lmodern}
\usepackage{graphicx}
\usepackage{array}
\usepackage{longtable}
\usepackage{booktabs}
\usepackage{tabularx}
\usepackage{enumitem}
\usepackage{xcolor}
\usepackage{titlesec}
\usepackage{fancyhdr}
\usepackage{tikz}
\usetikzlibrary{shapes.geometric, shapes.multipart, arrows.meta, positioning, fit, backgrounds, calc}
\usepackage[hidelinks]{hyperref}
\usepackage{caption}
\usepackage{placeins}

% ====================== KOLORY ======================
\definecolor{primary}{RGB}{30,90,150}
\definecolor{accent}{RGB}{0,140,120}
\definecolor{lightgray}{RGB}{240,242,245}
\definecolor{midgray}{RGB}{170,178,186}
\definecolor{ucfill}{RGB}{225,238,248}
\definecolor{compfill}{RGB}{235,245,238}

% ====================== STYLE NAGŁÓWKÓW ======================
\titleformat{\section}{\normalfont\Large\bfseries\color{primary}}{\thesection.}{0.6em}{}
\titleformat{\subsection}{\normalfont\large\bfseries\color{accent}}{\thesubsection}{0.6em}{}
\titlespacing*{\section}{0pt}{18pt}{8pt}

% ====================== STOPKA / NAGŁÓWEK ======================
\pagestyle{fancy}
\fancyhf{}
\renewcommand{\headrulewidth}{0.4pt}
\renewcommand{\footrulewidth}{0.4pt}
\fancyhead[L]{\small\color{primary}\textbf{RehabSense}}
\fancyhead[R]{\small Dokument Architektury}
\fancyfoot[L]{\small Politechnika Łódzka}
\fancyfoot[R]{\small Strona \thepage}

% ====================== POMOCNICZE ======================
\newcolumntype{Y}{>{\raggedright\arraybackslash}X}
\newcommand{\sysname}{\textbf{RehabSense}}
\setlist{itemsep=2pt, topsep=3pt}

% ---- Style diagramów ----
\tikzset{
  layer/.style={rectangle, rounded corners, draw=primary, thick, fill=lightgray,
                minimum width=3.4cm, minimum height=1cm, align=center, font=\small},
  ext/.style={rectangle, rounded corners, draw=accent, thick, fill=white,
              minimum width=2.8cm, minimum height=0.9cm, align=center, font=\small},
  arr/.style={-{Stealth[length=3mm]}, thick, draw=primary!70},
  uc/.style={ellipse, draw=primary, thick, fill=ucfill, align=center,
             font=\scriptsize, minimum width=2.9cm, minimum height=0.95cm, inner sep=2pt},
  comp/.style={rectangle, draw=accent, thick, fill=compfill, rounded corners=2pt,
               align=center, font=\scriptsize\bfseries, minimum width=2.7cm, minimum height=1cm},
  cls/.style={draw=primary, thick, fill=lightgray, rectangle split,
              rectangle split parts=2, font=\scriptsize, text width=3.0cm,
              align=left, rounded corners=2pt},
  part/.style={draw=primary, thick, fill=lightgray, rounded corners,
               minimum width=2.5cm, minimum height=0.8cm, font=\footnotesize\bfseries, align=center},
}
% Aktor (stick figure) jako pic
\tikzset{actorpic/.pic={
  \draw[thick] (0,0) circle (3.2mm);
  \draw[thick] (0,-3.2mm) -- (0,-12mm);
  \draw[thick] (-5mm,-7mm) -- (5mm,-7mm);
  \draw[thick] (0,-12mm) -- (-5mm,-20mm);
  \draw[thick] (0,-12mm) -- (5mm,-20mm);
}}
% Wiadomości w diagramach sekwencji
\newcommand{\msg}[4]{\draw[-{Stealth[length=2mm]},thick,primary] (#1 |- 0,#3) -- node[above,font=\scriptsize,midway,text=black]{#4} (#2 |- 0,#3);}
\newcommand{\rsg}[4]{\draw[-{Stealth[length=2mm]},dashed,accent] (#1 |- 0,#3) -- node[above,font=\scriptsize,midway,text=black]{#4} (#2 |- 0,#3);}

\begin{document}

% ====================== STRONA TYTUŁOWA ======================
\begin{titlepage}
\centering
\vspace*{2cm}
{\Huge\bfseries\color{primary} Dokument Architektury\par}
\vspace{0.5cm}
{\LARGE System wspomagania rehabilitacji\par}
\vspace{0.3cm}
{\Huge\bfseries RehabSense\par}
\vspace{1.2cm}
\rule{0.7\textwidth}{1.2pt}\par
\vspace{1.2cm}
{\large\textbf{Instytucja:} Politechnika Łódzka\par}
\vspace{0.4cm}
{\large\textbf{Dokument:} Dokumentacja\par}
\vspace{0.4cm}
{\large\textbf{Data:} 19 czerwca 2026\par}
\end{titlepage}

% ====================== LISTA CZŁONKÓW ZESPOŁU ======================
\section*{Lista członków zespołu}
\addcontentsline{toc}{section}{Lista członków zespołu}
\begin{table}[ht]
\centering
\begin{tabularx}{\textwidth}{>{\bfseries}p{4.2cm} p{3.5cm} Y}
\toprule
\textbf{Imię i nazwisko} & \textbf{Rola} & \textbf{Obszar} \\
\midrule
Kacper Błaszczyk & Project Manager & Koordynacja, architektura, wymagania. \\
Klim Hudzenko & Backend Developer & Backend (FastAPI), baza danych, API. \\
Maciej Miazek & Frontend Developer & Frontend (Next.js), moduł AI. \\
Remigiusz Tomecki & Frontend Developer & Frontend (Next.js), interfejs użytkownika. \\
Kacper Romek & Tester & Testy, zapewnienie jakości. \\
\bottomrule
\end{tabularx}
\caption{Skład zespołu projektowego RehabSense.}
\end{table}
\newpage

% ====================== SPIS TREŚCI ======================
\tableofcontents
\newpage

% ====================================================
\section{Wprowadzenie}
% ====================================================
\subsection{Cel dokumentu}
Niniejszy dokument przedstawia architekturę systemu \sysname{}. Opisuje przeznaczenie i funkcje oprogramowania, kluczowe przypadki użycia (wraz ze scenariuszami), strukturę komponentów, model danych (diagram klas) oraz interakcje między elementami systemu (diagramy sekwencji). Celem dokumentu jest dostarczenie spójnego obrazu rozwiązania, stanowiącego podstawę implementacji, testów i dalszego rozwoju.

\subsection{Zakres projektu}
\sysname{} to aplikacja webowa wspomagająca rehabilitację, wykorzystująca analizę wizyjną postawy ciała do oceny poprawności ćwiczeń wykonywanych przez pacjenta oraz wspierająca współpracę pacjenta z fizjoterapeutą. Projekt obejmuje warstwę prezentacji (frontend), logikę biznesową (backend / REST API), moduł analizy pozy (AI po stronie klienta) oraz warstwę danych (hostowana baza NeonDB).

\subsection{Odbiorcy projektu}
\begin{itemize}
  \item \textbf{Pacjenci} — osoby wymagające rehabilitacji
  \item \textbf{Fizjoterapeuci} — osoby przeprowadzające rehabilitację
  \item \textbf{Deweloperzy} — twórcy wykorzystujący funkcjonalności naszego systemu
\end{itemize}

% ====================================================
\section{Opis testowanego systemu}
% ====================================================
\subsection{Przeznaczenie oprogramowania}
\sysname{} to inteligentny interfejs człowiek–maszyna wspomagający diagnostykę, terapię i rehabilitację. System umożliwia wykonywanie ćwiczeń rehabilitacyjnych w domu pod nadzorem algorytmów AI, które w czasie rzeczywistym oceniają poprawność ruchu, a także wspiera komunikację i monitorowanie postępów we współpracy pacjenta z fizjoterapeutą.

\subsection{Główne funkcje aktorów systemu}
\begin{itemize}
  \item \textbf{Pacjent:} rejestracja/logowanie, parowanie z fizjoterapeutą, sesje ćwiczeń z analizą AI, panel statystyk i postępów, system passy, czat.
  \item \textbf{Fizjoterapeuta:} zarządzanie pacjentami, tworzenie i przypisywanie spersonalizowanych planów ćwiczeń, monitoring postępów, certyfikaty.
  \item \textbf{Administrator:} zarządzanie bazą zweryfikowanych fizjoterapeutów i kontami.
  \item \textbf{System (AI):} porównywanie ruchu użytkownika ze wzorcem ćwiczenia (MediaPipe + DTW).
\end{itemize}

\subsection{Architektura systemu}
System zaprojektowano w modelu trójwarstwowym (klient–serwer). Frontend (Next.js) komunikuje się z backendem (FastAPI) poprzez REST API zabezpieczone tokenami JWT. Analiza pozy odbywa się po stronie klienta (MediaPipe + DTW), a dane zapisywane są w hostowanej bazie NeonDB (PostgreSQL).

\begin{figure}[ht]
\centering
\begin{tikzpicture}[node distance=0.9cm and 1.4cm]
\node[layer] (front) {\textbf{Warstwa Prezentacji}\\ Frontend (Next.js)\\ \footnotesize Pacjent / Fizjoterapeuta / Admin};
\node[ext, below left=1.3cm and -0.5cm of front] (ai) {\textbf{Moduł AI}\\ MediaPipe + DTW\\ \footnotesize (analiza pozy w przeglądarce)};
\node[layer, below=2.6cm of front] (back) {\textbf{Warstwa Logiki}\\ Backend (FastAPI / Python)\\ \footnotesize Auth, Pacjent, Fizjo, Admin, Chat, Passy};
\node[layer, below=1.2cm of back] (db) {\textbf{Warstwa Danych}\\ NeonDB (PostgreSQL) + SQLAlchemy\\ \footnotesize profile, wyniki, plany, wiadomości};
\node[ext, right=2.2cm of back] (cloud) {\textbf{Usługi zewn.}\\ Cloudinary\\ \footnotesize (multimedia)};
\node[ext, right=2.2cm of front] (extapi) {\textbf{API zewn.}\\ \texttt{/external/v1}\\ \footnotesize (klucz API)};
\draw[arr] (front) -- node[right, font=\scriptsize] {REST / JWT} (back);
\draw[arr] (back) -- (db);
\draw[arr] (ai.south) |- (back.west);
\draw[arr] (back.east) -- (cloud.west);
\draw[arr] (extapi.south) -- (back.north east);
\end{tikzpicture}
\caption{Architektura trójwarstwowa systemu RehabSense.}
\end{figure}

\subsection{Integracje z systemami zewnętrznymi}
\begin{itemize}
  \item \textbf{NeonDB} — hostowana baza PostgreSQL jako usługa chmurowa.
  \item \textbf{Cloudinary} — przechowywanie i serwowanie materiałów wideo (instruktaże ćwiczeń, certyfikaty).
  \item \textbf{API zewnętrzne} (\texttt{/external/v1}) — udostępnianie danych systemom zewnętrznym, autoryzacja kluczem API.
  \item \textbf{Google MediaPipe} — biblioteka wizyjna (tasks-vision) ładowana po stronie klienta.
\end{itemize}

\newpage
% ====================================================
\section{Diagramy przypadków użycia}
% ====================================================
Poniższy diagram przedstawia główne przypadki użycia systemu \sysname{} oraz powiązanych z nimi aktorów.

\begin{figure}[ht]
\centering
\resizebox{0.98\textwidth}{!}{%
\begin{tikzpicture}

% --- Definicja nowej ikony dla Systemu (sześcian zamiast głowy) ---
\tikzset{
    systempic/.pic={
        % Głowa (sześcian izometryczny)
        \draw[thick] (0,0) -- (-0.25, 0.15) -- (-0.25, 0.45) -- (0, 0.6) -- (0.25, 0.45) -- (0.25, 0.15) -- cycle; 
        \draw[thick] (0,0) -- (0, 0.3); 
        \draw[thick] (-0.25, 0.45) -- (0, 0.3) -- (0.25, 0.45); 
        % Tułów
        \draw[thick] (0,0) -- (0,-0.6);
        % Ramiona (przecinające się u podstawy sześcianu)
        \draw[thick] (-0.4,0) -- (0.4,0);
        % Nogi
        \draw[thick] (-0.3,-1) -- (0,-0.6) -- (0.3,-1);
    }
}

% --- Granica systemu (standardowa zakładka UML) ---
% Zwiększono szerokość ramki z 8.5 na 9.5 z prawej strony
\draw[thick] (-5.5, -8.0) rectangle (10, 7.5);
\draw[thick, fill=white] (-5.5, 7.5) -- (-5.5, 8.3) -- (-2.5, 8.3) -- (-2.2, 7.5) -- cycle;
\node[font=\bfseries, right] at (-5.3, 7.9) {RehabSense};

% --- Przypadki użycia (wewnątrz) ---
\node[uc] (passy)        at (0.0, 6.5)   {Wyświetlanie \\ passy};
\node[uc] (sesja)        at (1.5, 4.5)   {Wykonywanie sesji \\ ćwiczeń};
% Przesunięto X z 6.0 na 7.5, by wydłużyć strzałki include
\node[uc] (sledzenie)    at (7.5, 5.5)   {Śledzenie poprawności \\ wykonywanych ćwiczeń};
\node[uc] (aktualizacja) at (7.5, 3.0)   {Aktualizacja postępu \\ rehabilitacji};
\node[uc] (plan)         at (2.5, 1.5)   {Wyświetlanie planu \\ rehabilitacji};
\node[uc] (polaczenie)   at (-1.0, 2.5)  {Połączenie z \\ fizjoterapeutą};
\node[uc] (statystyki)   at (-3.0, 0.5)  {Wyświetlanie \\ statystyk \\ rehabilitacji};
\node[uc] (komunikacja)  at (1.0, -1.0)  {Komunikacja z \\ pacjentem/ \\ fizjoterapeutą};
\node[uc] (tworzenie)    at (2.5, -3.0)  {Tworzenie planów \\ ćwiczeń};
\node[uc] (lista)        at (-2.5, -4.5) {Przeglądanie listy \\ pacjentów};
\node[uc] (admin_uc)     at (0.0, -6.5)  {Dodawanie profili \\ fizjoterapeutów};

% --- Aktorzy ---
\pic at (-7.5, 3.5) {actorpic};
\node[font=\small\bfseries] at (-7.5, 1.3) {Pacjent}; 

\pic at (-7.5, -2.0) {actorpic};
\node[font=\small\bfseries] at (-7.5, -4.2) {Fizjoterapeuta}; 

\pic at (-7.5, -7.0) {actorpic};
\node[font=\small\bfseries] at (-7.5, -9.2) {Administrator}; 

% Przesunięto System z 10.0 na 11.5, by dopasować się do nowej szerokości ramki
\pic at (11.5, 4.25) {systempic};
\node[font=\small\bfseries] at (11.5, 2.45) {System}; 

% --- Powiązania Pacjent ---
\draw (-7.1, 3.5) -- (passy.west);
\draw (-7.1, 3.5) -- (sesja.west);
\draw (-7.1, 3.5) -- (plan.west);
\draw (-7.1, 3.5) -- (polaczenie.west);
\draw (-7.1, 3.5) -- (statystyki.west);
\draw (-7.1, 3.5) -- (komunikacja.west);

% --- Powiązania Fizjoterapeuta ---
\draw (-7.1, -2.0) -- (statystyki.west);
\draw (-7.1, -2.0) -- (komunikacja.west);
\draw (-7.1, -2.0) -- (tworzenie.west);
\draw (-7.1, -2.0) -- (lista.west);

% --- Powiązania Administrator ---
\draw (-7.1, -7.0) -- (admin_uc.west);

% --- Powiązania System ---
\draw (11.1, 4.25) -- (sledzenie.east);
\draw (11.1, 4.25) -- (aktualizacja.east);

% --- Zależności <<include>> ---
\draw[-latex, dashed] (sesja.east) -- node[midway, sloped, above, font=\scriptsize] {<<include>>} (sledzenie.west);
\draw[-latex, dashed] (sesja.east) -- node[midway, sloped, above, font=\scriptsize] {<<include>>} (aktualizacja.west);

\end{tikzpicture}}
\caption{Diagram przypadków użycia systemu RehabSense.}
\end{figure}

\FloatBarrier
\clearpage
\subsection{Scenariusze przypadków użycia}
Poniżej opisano scenariusze trzech ważnych przypadków użycia.

% ---------- 1. ----------
\begin{table}[ht]
\small
\renewcommand{\arraystretch}{1.2}
\begin{tabularx}{\textwidth}{|>{\bfseries}p{0.22\textwidth}|Y|}
\hline
Nazwa & Parowanie pacjenta i fizjoterapeuty \\ \hline
Poziom ważności & Wysoki \\ \hline
Aktorzy & Pacjent, Fizjoterapeuta \\ \hline
Krótki opis & Sparowanie pacjenta z fizjoterapeutą. \\ \hline
Warunki wstępne & Pacjent i Fizjoterapeuta są zarejestrowani; pacjent jest zalogowany. \\ \hline
Warunki końcowe & Profile pacjenta i fizjoterapeuty zostają połączone (status \texttt{ZAAKCEPTOWANE}). \\ \hline
Główny przepływ & 1. Pacjent wyszukuje fizjoterapeutę w panelu wyboru fizjoterapeuty.\newline
2. Pacjent wybiera fizjoterapeutę i wysyła prośbę o połączenie.\newline
3. Fizjoterapeuta zatwierdza prośbę.\newline
4. System wysyła na chat między pacjentem, a fizjoterapeutą komunikat o poprawnym sparowaniu \\ \hline
Przepływy alternatywne & \textbf{2a.} Pacjent wyświetla certyfikat fizjoterapeuty. \newline \textbf{3a.} Fizjoterapeuta odrzuci prośbę. (powrót do kroku 2). \\ \hline
Wymagania specjalne & brak \\ \hline
\end{tabularx}
\caption{Scenariusz 1 — Parowanie pacjenta i fizjoterapeuty.}
\end{table}
\FloatBarrier

% ---------- 2.  ----------
\begin{table}[ht]
\small
\renewcommand{\arraystretch}{1.2}
\begin{tabularx}{\textwidth}{|>{\bfseries}p{0.22\textwidth}|Y|}
\hline
Nazwa & Sesja rehabilitacji (analiza AI) \\ \hline
Poziom ważności & Krytyczny \\ \hline
Aktorzy & Pacjent, System \\ \hline
Krótki opis & Pacjent wykonuje ćwiczenia przed kamerą z analizą postawy ciała i postępem passy. \\ \hline
Warunki wstępne & Pacjent jest zalogowany; ma aktywny plan rehabilitacji i kamerke internetowa. \\ \hline
Warunki końcowe & Wyniki sesji (powtórzenia, dokładność) zostają zapisane. Passa i ewentualne odznaki zaktualizowane. \\ \hline
Główny przepływ & 1. Pacjent rozpoczyna sesję; system pobiera aktywny plan i ćwiczenia.\newline
2. Dla ćwiczenia system pobiera dane wzorcowe.\newline
3. System wykrywa ruch pacjenta i algorytm DTW porównuje szkielet MediaPipe ze wzorcem.\newline
4. System wyświetla informację zwrotną w czasie rzeczywistym (dokładność).\newline
5. System zapisuje sesję i wyniki ćwiczeń.\newline
6. System nalicza passę i przyznaje odznaki. \\ \hline
Przepływy alternatywne & \textbf{3a.} Niska jakość obrazu / brak wykrycia postawy — system informuje o problemie i prosi o korektę warunków.\newline
\textbf{5a.} Błąd zapisu wyników — sesja zostaje oznaczona jako nieukończona. \\ \hline
Wymagania specjalne & Analiza w czasie rzeczywistym (ok. 30 FPS). \\ \hline
\end{tabularx}
\caption{Scenariusz 2 — Sesja rehabilitacji (analiza AI).}
\end{table}
\FloatBarrier

% ---------- UC3 ----------
\begin{table}[ht]
\small
\renewcommand{\arraystretch}{1.25}
\begin{tabularx}{\textwidth}{|>{\bfseries}p{0.22\textwidth}|Y|}
\hline
Nazwa & Tworzenie i przypisanie planu rehabilitacji \\ \hline
Numer & 3 \\ \hline
Poziom ważności & Wysoki \\ \hline
Aktorzy & Fizjoterapeuta, System \\ \hline
Krótki opis & Fizjoterapeuta tworzy spersonalizowany plan ćwiczeń i przypisuje go wybranemu pacjentowi, który otrzymuje powiadomienie. \\ \hline
Warunki wstępne & Fizjoterapeuta jest zalogowany i ma co najmniej jednego przypisanego pacjenta oraz dostępne są ćwiczenia (globalne lub własne). \\ \hline
Warunki końcowe & Nowy aktywny plan zostaje utworzony. Poprzedni aktywny plan pacjenta zostaje dezaktywowany. Pacjent otrzymuje powiadomienie systemowe. \\ \hline
Główny przepływ & 1. Fizjoterapeuta otwiera kreator planu i pobiera listę swoich pacjentów oraz ćwiczeń.\newline
2. Wybiera pacjenta oraz ćwiczenia wraz z liczbą serii i powtórzeń.\newline
3. System dezaktywuje poprzedni aktywny plan pacjenta.\newline
4. System tworzy nowy plan i przypisuje do niego ćwiczenia.\newline
5. System wysyła pacjentowi powiadomienie o aktualizacji planu. \\ \hline
Przepływy alternatywne & \textbf{2a.} Wybrane ćwiczenie nie posiada przypisanego wideo — system odrzuca utworzenie planu i informuje o błędzie. \\ \hline
Wymagania specjalne & Każde ćwiczenie musi mieć filmik instruktażowy. \\ \hline
\end{tabularx}
\caption{Scenariusz 3 — Tworzenie i przypisanie planu rehabilitacji.}
\end{table}
\FloatBarrier

% ====================================================
\FloatBarrier
\section{Diagram komponentów}
% ====================================================
Poniższy diagram przedstawia logiczny podział systemu na komponenty oraz ich zależności.

\begin{figure}[ht]
\centering
\resizebox{\textwidth}{!}{%
\begin{tikzpicture}[node distance=0.8cm and 3.5cm]

% Definicja ikony komponentu (UML)
\newcommand{\drawcompicon}[1]{
  \draw[accent, thick, fill=white] ($(#1.north east) + (-0.5,-0.2)$) rectangle ($(#1.north east) + (-0.1,-0.6)$);
  \draw[accent, thick, fill=white] ($(#1.north east) + (-0.6,-0.3)$) rectangle ($(#1.north east) + (-0.4,-0.35)$);
  \draw[accent, thick, fill=white] ($(#1.north east) + (-0.6,-0.45)$) rectangle ($(#1.north east) + (-0.4,-0.5)$);
}

% Węzły - Środek
\node[comp, minimum width=4.5cm, minimum height=1.5cm] (mp) {Moduł Pacjenta};
\node[comp, minimum width=4.5cm, minimum height=1.5cm, above=of mp] (ma) {Moduł Administracji};
\node[comp, minimum width=4.5cm, minimum height=1.5cm, below=of mp] (mf) {Moduł Fizjoterapeuty};
\node[comp, minimum width=4.5cm, minimum height=1.5cm, below=of mf] (mai) {Moduł AI};

% Rysowanie ikon UML dla środka
\drawcompicon{mp}
\drawcompicon{ma}
\drawcompicon{mf}
\drawcompicon{mai}

% Lewa strona - Interfejs Graficzny
\node[comp, minimum width=4.5cm, minimum height=2cm, left=3.5cm of mp, yshift=-1.15cm] (ig) {Interfejs Graficzny};
\drawcompicon{ig}

% Prawa strona - Baza Danych
\node[comp, minimum width=4.5cm, minimum height=2cm, right=3.5cm of mp, yshift=-1.15cm] (bd) {Baza Danych};
\drawcompicon{bd}

% Konfiguracja symbolu interfejsu (łapka i kółko)
\tikzset{
    interface/.pic={
        % Wypełnienie tła, by ukryć fragment linii wewnątrz łapki
        \fill[white] (0, 2.5mm) arc (90:270:2.5mm) -- cycle;
        % Łapka (gniazdo) po lewej stronie
        \draw[thick, draw=accent] (0, 2.5mm) arc (90:270:2.5mm);
        % Kółko (wtyczka) na środku
        \node[circle, draw=accent, fill=white, inner sep=0pt, minimum size=3.5mm] at (0,0) {};
    }
}

% Interfejsy - Lewa strona (IG do Modułów)
\draw[thick, draw=accent] (ig.east) -- pic[pos=0.5, sloped] {interface} node[pos=0.5, above left, yshift=1mm, font=\scriptsize] {IAdmin} (ma.west);

\draw[thick, draw=accent] (ig.east) -- pic[pos=0.5, sloped] {interface} node[pos=0.5, above left, yshift=1mm, font=\scriptsize] {IPatient} (mp.west);

\draw[thick, draw=accent] (ig.east) -- pic[pos=0.5, sloped] {interface} node[pos=0.5, below left, yshift=-1mm, font=\scriptsize] {IPhysiotherapist} (mf.west);

\draw[thick, draw=accent] (ig.east) -- pic[pos=0.5, sloped] {interface} node[pos=0.5, below left, yshift=-1mm, font=\scriptsize] {IAIAnalysis} (mai.west);

% Interfejsy - Prawa strona (Moduły do Bazy Danych)
\draw[thick, draw=accent] (ma.east) -- pic[pos=0.5, sloped] {interface} node[pos=0.5, above right, yshift=1mm, font=\scriptsize] {IDatabase} (bd.west);

\draw[thick, draw=accent] (mp.east) -- pic[pos=0.5, sloped] {interface} node[pos=0.5, above right, yshift=1mm, font=\scriptsize] {IDatabase} (bd.west);

\draw[thick, draw=accent] (mf.east) -- pic[pos=0.5, sloped] {interface} node[pos=0.5, below right, yshift=-1mm, font=\scriptsize] {IDatabase} (bd.west);

\draw[thick, draw=accent] (mai.east) -- pic[pos=0.5, sloped] {interface} node[pos=0.5, below right, yshift=-1mm, font=\scriptsize] {IDatabase} (bd.west);

\end{tikzpicture}}
\caption{Diagram komponentów systemu RehabSense.}
\end{figure}

\subsection{Opis komponentów}
\begin{itemize}
  \item \textbf{Interfejs UI (Next.js/React):} warstwa prezentacji — strony pulpitu, formularze, panele statystyk; komunikuje się z backendem przez REST API.
  \item \textbf{Moduł AI (PoseDetector + DTW):} działa w przeglądarce; wykorzystuje MediaPipe do detekcji pozy oraz DTW do porównania ruchu ze wzorcem i obliczenia dokładności.
  \item \textbf{Auth (JWT, RoleChecker):} uwierzytelnianie tokenami JWT i autoryzacja oparta na rolach (pacjent / fizjoterapeuta / admin).
  \item \textbf{API biznesowe:} routery FastAPI dla modułów Pacjenta, Fizjoterapeuty i Administratora (plany, sesje, pacjenci, czat).
  \item \textbf{Logika domenowa:} reguły biznesowe — parowanie, passy, odznaki (\texttt{badges}).
  \item \textbf{AI API:} udostępnia dane wzorcowe ćwiczeń dla modułu AI klienta.
  \item \textbf{External API:} integracja z systemami zewnętrznymi (autoryzacja kluczem API).
  \item \textbf{NeonDB:} warstwa danych — trwałe przechowywanie profili, planów, sesji, wyników i wiadomości.
  \item \textbf{Cloudinary:} hostowanie materiałów multimedialnych (wideo instruktażowe, certyfikaty).
\end{itemize}

% ====================================================
\FloatBarrier
\section{Diagram klas}
% ====================================================
Diagram przedstawia kluczowe encje modelu danych (mapowane przez SQLAlchemy) wraz z najważniejszymi atrybutami i powiązaniami.

\begin{figure}[ht]
\centering
\resizebox{\textwidth}{!}{%
\begin{tikzpicture}[every node/.style={font=\scriptsize}, cls/.append style={text width=5.0cm}]

% User
\node[cls] (user) at (-15,0)
  {\textbf{User}\nodepart{second} +user\_id: Integer [PK] \\ +first\_name: String \\ +last\_name: String \\ +email: String \\ +password: String \\ +role: String \\ +api\_key: String};

% Message
\node[cls] (msg) at (-15,-5)
  {\textbf{Message}\nodepart{second} +id: Integer [PK] \\ +sender\_id: Integer [FK] \\ +receiver\_id: Integer [FK] \\ +content: String \\ +timestamp: DateTime \\ +is\_read: Boolean};

% Patient
\node[cls] (patient) at (-7,0)
  {\textbf{Patient}\nodepart{second} +user\_id: Integer [PK, FK] \\ +streak\_count: Integer \\ +last\_activity: Date};

% Physiotherapist
\node[cls] (physio) at (-7,-5)
  {\textbf{Physiotherapist}\nodepart{second} +user\_id: Integer [PK, FK] \\ +specialization: String \\ +is\_verified: Boolean \\ +is\_available: Boolean};

% ProgressNote
\node[cls] (note) at (-7,-10)
  {\textbf{ProgressNote}\nodepart{second} +note\_id: Integer [PK] \\ +patient\_id: Integer [FK] \\ +physio\_id: Integer [FK] \\ +created\_at: DateTime \\ +note\_content: String \\ +pain\_level: Integer \\ +mobility\_level: Integer};

% PatientPhysiotherapist
\node[cls] (pp) at (1,-5)
  {\textbf{PatientPhysiotherapist}\nodepart{second} +id: Integer [PK] \\ +patient\_id: Integer [FK] \\ +physio\_id: Integer [FK] \\ +status: String};

% RehabPlan
\node[cls] (plan) at (9,0)
  {\textbf{RehabPlan}\nodepart{second} +rehab\_id: Integer [PK] \\ +physio\_id: Integer [FK] \\ +patient\_id: Integer [FK] \\ +title: String \\ +is\_active: Boolean};

% RehabPlanExercise
\node[cls] (rpe) at (9,-10)
  {\textbf{RehabPlanExercise}\nodepart{second} +id: Integer [PK] \\ +rehab\_id: Integer [FK] \\ +exercise\_id: Integer [FK] \\ +reps\_nr: Integer \\ +sets\_nr: Integer};

% Session
\node[cls] (sess) at (16,0)
  {\textbf{Session}\nodepart{second} +session\_id: Integer [PK] \\ +rehab\_id: Integer [FK] \\ +patient\_id: Integer [FK] \\ +title: String \\ +is\_active: Boolean \\ +created\_at: DateTime};

% ExerciseResult
\node[cls] (res) at (16,-5)
  {\textbf{ExerciseResult}\nodepart{second} +result\_id: Integer [PK] \\ +exercise\_id: Integer [FK] \\ +session\_id: Integer [FK] \\ +reps\_completed: Integer \\ +avg\_accuracy: Float \\ +ai\_feedback: String \\ +max\_rom: Float \\ +pain\_level: Integer};

% Exercise
\node[cls] (ex) at (16,-10)
  {\textbf{Exercise}\nodepart{second} +exercise\_id: Integer [PK] \\ +name: String \\ +description: Text \\ +video\_url: String \\ +mediapipe\_pattern\_data: Text};

% Relacje
\draw[-] (user.east) -- node[above, font=\tiny]{has profile} node[pos=0.05, above, font=\tiny]{1} node[pos=0.95, above, font=\tiny]{0..1} (patient.west);

\draw[-] (user.south east) -- node[above, sloped, font=\tiny]{has profile} node[pos=0.05, above, font=\tiny]{1} node[pos=0.95, above, font=\tiny]{0..1} (physio.north west);

\draw[-] (user.south) -- node[right, font=\tiny]{sends / receives} node[pos=0.05, right, font=\tiny]{1} node[pos=0.95, right, font=\tiny]{*} (msg.north);

\draw[-] (patient.south west) -- ++(-0.6, 0) |- node[pos=0.25, right, font=\tiny]{has} node[pos=0.05, above, font=\tiny]{1} node[pos=0.95, above, font=\tiny]{*} (note.west);

\draw[-] (physio.south) -- node[right, font=\tiny]{writes} node[pos=0.05, right, font=\tiny]{1} node[pos=0.95, right, font=\tiny]{*} (note.north);

\draw[-] (patient.south east) -- node[above, sloped, font=\tiny]{is associated with} node[pos=0.05, above, font=\tiny]{1} node[pos=0.95, above, font=\tiny]{*} (pp.north west);

\draw[-] (physio.east) -- node[above, font=\tiny]{is associated with} node[pos=0.05, above, font=\tiny]{1} node[pos=0.95, above, font=\tiny]{*} (pp.west);

\draw[-] (patient.east) -- node[above, font=\tiny]{assigned to} node[pos=0.05, above, font=\tiny]{1} node[pos=0.95, above, font=\tiny]{*} (plan.west);

\draw[-] (physio.north east) -- node[above, sloped, font=\tiny]{creates} node[pos=0.05, above, font=\tiny]{1} node[pos=0.95, above, font=\tiny]{*} (plan.south west);

\draw[-] (plan.south) -- node[right, font=\tiny]{contains} node[pos=0.05, right, font=\tiny]{1} node[pos=0.95, right, font=\tiny]{*} (rpe.north);

\draw[-] (plan.east) -- node[above, font=\tiny]{has} node[pos=0.05, above, font=\tiny]{1} node[pos=0.95, above, font=\tiny]{*} (sess.west);

\draw[-] (patient.north) -- ++(0, 2.0) -| node[pos=0.25, above, font=\tiny]{participates in} node[pos=0.02, right, font=\tiny]{1} node[pos=0.98, right, font=\tiny]{*} (sess.north);

\draw[-] (rpe.east) -- node[above, font=\tiny]{part of} node[pos=0.05, above, font=\tiny]{*} node[pos=0.95, above, font=\tiny]{1} (ex.west);

\draw[-] (sess.south) -- node[right, font=\tiny]{records} node[pos=0.05, right, font=\tiny]{1} node[pos=0.95, right, font=\tiny]{*} (res.north);

\draw[-] (res.south) -- node[right, font=\tiny]{result of} node[pos=0.05, right, font=\tiny]{*} node[pos=0.95, right, font=\tiny]{1} (ex.north);
\end{tikzpicture}}
\caption{Diagram klas (model danych) systemu RehabSense.}
\end{figure}

% ====================================================
\FloatBarrier
\clearpage
\section{Diagramy interakcji}
% ====================================================
Poniżej przedstawiono diagramy sekwencji dla trzech opisanych wcześniej przypadków użycia.

% ---------- SD1: Parowanie ----------
\subsection{Diagram sekwencji do scenariusza 1 - Parowanie pacjenta i fizjoterapeuty}
\begin{figure}[ht]
\centering
\resizebox{\textwidth}{!}{%
\begin{tikzpicture}
\node[part] (pac) at (0,0) {Pacjent};
\node[part] (fe)  at (3.6,0) {Frontend};
\node[part] (be)  at (7.2,0) {Backend API};
\node[part] (db)  at (10.8,0) {NeonDB};
\node[part] (fiz) at (14.4,0) {Fizjoterapeuta};
\foreach \p in {pac,fe,be,db,fiz} \draw[dashed,midgray] (\p.south) -- (\p |- 0,-11.2);
\msg{pac}{fe}{-1.0}{wybór „Szukaj fizjo'' + problem}
\msg{fe}{be}{-1.8}{POST /pairing/request}
\msg{be}{db}{-2.6}{find\_available\_physio()}
\rsg{db}{be}{-3.4}{dane fizjoterapeuty}
\msg{be}{db}{-4.2}{create\_pairing\_request (OCZEKUJACE)}
\rsg{be}{fe}{-5.0}{200 (request\_id)}
\msg{fiz}{fe}{-5.8}{akceptacja zapytania w panelu}
\msg{fe}{be}{-6.6}{POST /pairing/\{request\_id\}/physio-respond}
\msg{be}{db}{-7.4}{PatientPhysiotherapist = ZAAKCEPTOWANE}
\msg{be}{db}{-8.2}{Message [SYSTEM:CONNECT]}
\rsg{be}{pac}{-9.0}{komunikat na czacie}
\end{tikzpicture}}
\caption{Diagram sekwencji 1 Parowanie pacjenta i fizjoterapeuty.}
\end{figure}

% ---------- SD2: Sesja ----------
\FloatBarrier
\subsection{Diagram sekwencji do scenariusza 2 - Sesja rehabilitacji (analiza AI)}
\begin{figure}[ht]
\centering
\resizebox{\textwidth}{!}{%
\begin{tikzpicture}
\node[part] (pac) at (0,0) {Pacjent};
\node[part] (fe)  at (3.6,0) {Frontend + AI};
\node[part] (be)  at (7.6,0) {Backend API};
\node[part] (db)  at (11.4,0) {NeonDB};
\foreach \p in {pac,fe,be,db} \draw[dashed,midgray] (\p.south) -- (\p |- 0,-10.2);
\msg{pac}{fe}{-1.0}{rozpoczyna sesję}
\msg{fe}{be}{-1.8}{GET /patient/my-plan}
\msg{be}{db}{-2.6}{query RehabPlan + ćwiczenia}
\rsg{be}{fe}{-3.4}{plan (lista ćwiczeń)}
\msg{fe}{be}{-4.2}{GET /ai/pattern/\{exercise\_id\}}
\rsg{be}{fe}{-5.0}{mediapipe\_pattern\_data}
\msg{pac}{fe}{-5.8}{wykonuje ruch (kamera)}
\node[draw=accent,fill=white,font=\scriptsize,align=center] at (fe |- 0,-6.7) {MediaPipe + DTW:\\ detekcja pozy, dokładność};
\rsg{fe}{pac}{-7.6}{feedback w czasie rzeczywistym}
\msg{fe}{be}{-8.4}{POST /patient/submit-session}
\msg{be}{db}{-9.2}{zapis Session + ExerciseResult, passa, odznaki}
\rsg{be}{fe}{-10.0}{streak, longest\_streak, new\_badges}
\end{tikzpicture}}
\caption{Diagram sekwencji 2. Sesja rehabilitacji z analizą AI.}
\end{figure}

% ---------- SD3: Tworzenie planu ----------
\FloatBarrier
\clearpage
\subsection{Diagram sekwencji do scenariusza 3 - Tworzenie i przypisanie planu rehabilitacji}
\begin{figure}[ht]
\centering
\resizebox{\textwidth}{!}{%
\begin{tikzpicture}
\node[part] (fiz) at (0,0) {Fizjoterapeuta};
\node[part] (fe)  at (3.6,0) {Frontend};
\node[part] (be)  at (7.2,0) {Backend API};
\node[part] (db)  at (10.8,0) {NeonDB};
\foreach \p in {fiz,fe,be,db} \draw[dashed,midgray] (\p.south) -- (\p |- 0,-9.4);
\msg{fiz}{fe}{-1.0}{otwiera kreator planu}
\msg{fe}{be}{-1.8}{GET /physio/my-patients}
\msg{be}{db}{-2.6}{query połączeń pacjentów}
\rsg{be}{fe}{-3.4}{lista pacjentów}
\msg{fe}{be}{-4.2}{GET /physio/exercises}
\rsg{be}{fe}{-5.0}{lista ćwiczeń}
\msg{fiz}{fe}{-5.8}{wybór pacjenta, ćwiczeń, serii/powt.}
\msg{fe}{be}{-6.6}{POST /physio/create-plan}
\msg{be}{db}{-7.4}{dezaktywacja starego + RehabPlan/RehabPlanExercise}
\msg{be}{db}{-8.2}{Message [SYSTEM:PLAN\_UPDATE]}
\rsg{be}{fe}{-9.0}{plan utworzony}
\end{tikzpicture}}
\caption{Diagram sekwencji 3. Tworzenie i przypisanie planu rehabilitacji.}
\end{figure}

% ====================================================
\FloatBarrier
\section{Stack technologiczny}
% ====================================================
\begin{longtable}{p{0.20\textwidth} p{0.24\textwidth} p{0.46\textwidth}}
\toprule
\textbf{Obszar} & \textbf{Technologia} & \textbf{Uzasadnienie wyboru} \\
\midrule
\endhead
Frontend & Next.js 16, React 19 & Dojrzały ekosystem, renderowanie hybrydowe (SSR/CSR), komponentowość i wydajność; naturalne wsparcie dla aplikacji webowej działającej w przeglądarce. \\
\addlinespace
Wygląd & Tailwind CSS & Szybkie budowanie spójnego, responsywnego interfejsu, utility-first ogranicza ilość własnego CSS. \\
\addlinespace
Stan / wykresy & Zustand, Chart.js & Globalny stan (Zustand) oraz gotowe, czytelne wizualizacje postępów (Chart.js). \\
\addlinespace
Analiza ruchu (AI) & Google MediaPipe + DTW & MediaPipe zapewnia dokładną detekcję pozy w czasie rzeczywistym po stronie klienta (bez wysyłania obrazu na serwer — prywatność); DTW pozwala porównać sekwencję ruchu ze wzorcem mimo różnic tempa. \\
\addlinespace
Backend & Python, FastAPI & Wysoka wydajność, automatyczna walidacja (Pydantic) i dokumentacja OpenAPI, szybki rozwój; Python ułatwia ewentualną pracę z danymi/AI. \\
\addlinespace
ORM & SQLAlchemy 2.0 & Nowoczesny ORM, niezależność od konkretnej bazy (SQLite lokalnie, PostgreSQL produkcyjnie), czytelne modele i relacje. \\
\addlinespace
Baza danych & NeonDB (PostgreSQL) & Hostowana baza PostgreSQL jako usługa chmurowa — brak zarządzania infrastrukturą i skalowalność. \\
\addlinespace
Bezpieczeństwo & PyJWT, bcrypt & Standard tokenów JWT do autoryzacji bezstanowej; bcrypt do bezpiecznego hashowania haseł. \\
\addlinespace
Multimedia & Cloudinary & Wyspecjalizowane hostowanie i transformacja wideo/obrazów, odciąża backend i bazę. \\
\addlinespace
Konteneryzacja & Docker / Docker Compose & Powtarzalne środowisko uruchomieniowe (frontend + backend), łatwe wdrożenie i testy zbliżone do produkcji. \\
\bottomrule
\end{longtable}

% ====================================================
\section{Podsumowanie}
% ====================================================
Architektura systemu \sysname{} opiera się na klarownym podziale trójwarstwowym, w którym warstwa prezentacji (Next.js) wraz z modułem AI działającym po stronie klienta odpowiada za interakcję i analizę ruchu, warstwa logiki (FastAPI) realizuje reguły biznesowe i bezpieczeństwo, a warstwa danych (NeonDB) zapewnia trwałe i skalowalne przechowywanie informacji.

Przedstawione przypadki użycia — \textbf{parowanie pacjenta z fizjoterapeutą}, \textbf{sesja rehabilitacji z analizą AI} oraz \textbf{tworzenie i przypisanie planu rehabilitacji} — obejmują rdzeń funkcjonalny systemu i zostały zilustrowane scenariuszami oraz diagramami sekwencji. Diagram komponentów oraz diagram klas pokazują spójność między logicznym podziałem odpowiedzialności a modelem danych.

Wybrany stack technologiczny łączy wydajność, bezpieczeństwo i prywatność (analiza obrazu lokalnie, bez przesyłania na serwer) oraz niski koszt utrzymania dzięki rozwiązaniom hostowanym (NeonDB, Cloudinary). Tak zaprojektowana architektura jest modularna i gotowa na dalszy rozwój — m.in. rozbudowę modułu AI, dodawanie nowych typów ćwiczeń oraz integracje przez udostępnione API zewnętrzne.

\end{document}
